% Sustentacion del analisis exploratorio de imagenes
% Compilar: latexmk -pdf -outdir=build/slides slides/presentacion.tex

\documentclass[11pt]{beamer}

\usetheme{Madrid}
\usecolortheme{beaver}
\setbeamertemplate{navigation symbols}{}
\setbeamertemplate{caption}[numbered]

\usepackage[utf8]{inputenc}
\usepackage[T1]{fontenc}
\usepackage[spanish]{babel}
\usepackage{lmodern}
\usepackage{graphicx}
\usepackage{booktabs}

\graphicspath{{slides/}{reports/}{figures/fase1/}{figures/fase2/}{figures/fase3/}{./}}

\makeatletter
\providecommand{\input@path}{}
\edef\input@path{{reports/tablas/}{./}}
\makeatother
\input{cifras}

\title[Estabilidad de descriptores]{Analisis exploratorio de imagenes y tabla de caracteristicas}
\subtitle{Estabilidad de los descriptores morfologicos del campo pulmonar}
\author[Trujillo \and Martinez]{Juanita Trujillo Narvaez \\ Karol Isabella Martinez Villarreal}
\institute[ECI]{Ingenieria Biomedica \\ PSIM - 2026II - 80 \\ Docente: Ing. Pablo Eduardo Caicedo-Rodriguez, Ph.D.}
\date{\today}

\logo{\includegraphics[height=0.7cm]{Escuela_Rosario_logo.png}}

\begin{document}

% ---------------------------------------------------------------- 1
\begin{frame}[plain]
    \titlepage
\end{frame}

% ---------------------------------------------------------------- 2
\begin{frame}{Problema, unidad y salida}
\footnotesize
\begin{columns}[T]
\begin{column}{0.52\textwidth}
\textbf{El problema}

La delimitacion del campo pulmonar precede a toda medicion y es inestable:
depende del detector, la exposicion y la patologia. Si la mascara se desplaza,
los descriptores heredan esa variabilidad.

\vspace{0.3cm}
\textbf{Objetivo}

Determinar que combinacion de realce y umbralizacion produce descriptores con
menor error frente a la referencia, \textbf{consistente entre sitios}.
\end{column}

\begin{column}{0.46\textwidth}
\begin{block}{Contrato de datos}
\scriptsize
\begin{tabular}{@{}ll@{}}
Unidad & imagen $\times$ estrategia \\
Clave & \texttt{sample\_id} \\
U. independiente & la imagen \\
Entrada & radiografia PA \\
Salida & mascara (espacial) \\
Referencia & trazada por personas \\
Metodo & clasico, sin etapa supervisada \\
Evaluacion & Dice y error relativo \\
\end{tabular}
\end{block}
\end{column}
\end{columns}

\vspace{0.2cm}
\scriptsize
La tabla acompana al metodo y no lo sustituye: la salida espacial se conserva
como mascara vinculada. No hay columna objetivo en sentido supervisado.
\end{frame}

% ---------------------------------------------------------------- 3
\begin{frame}{Datos, calidad y particion}
\footnotesize
\begin{columns}[T]
\begin{column}{0.45\textwidth}
\textbf{Cohorte}: \NImagenesTotal{} radiografias de dos sitios,
\NConMascara{} con mascara de referencia.

\vspace{0.25cm}
\textbf{Hallazgos de la auditoria}
\begin{itemize}\scriptsize
    \item Montgomery es de 8 bits, no 12 como indica la ficha
    \item Ninguna imagen tiene la escala invertida
    \item 335 de 800 estan en orientacion horizontal
    \item Las 9 imagenes mas pequenas son los 9 pacientes menores de dos anos
    \item 96 imagenes sin mascara: participan en la caracterizacion, no en la evaluacion
\end{itemize}

\vspace{0.2cm}
\textbf{Particion} 70/15/15 por unidad independiente, estratificada por origen
y clase, semilla fija. Los dos grupos de sujeto demostrados no se separan.
\end{column}

\begin{column}{0.53\textwidth}
\begin{figure}
\centering
\includegraphics[width=\textwidth]{dimensiones.png}
\end{figure}
\scriptsize
Montgomery entrega siempre el mismo tamano de archivo; Shenzhen lo ajusta al
paciente. El area sin senal alcanza el 67\% en el caso mas colimado.
\end{column}
\end{columns}
\end{frame}

% ---------------------------------------------------------------- 4
\begin{frame}{Hallazgos del analisis exploratorio}
\footnotesize
\begin{figure}
\centering
\includegraphics[width=0.92\textwidth]{intensidad.png}
\end{figure}

\vspace{-0.2cm}
\begin{columns}[T]
\begin{column}{0.48\textwidth}
\textbf{El sitio pesa mas que la clase}
\scriptsize

Mediana de intensidad: 111,0 en Montgomery frente a 173,5 en Shenzhen.
Dentro de Montgomery, 110,0 sin hallazgos y 111,0 con hallazgos.

Un descriptor que separase las clases estaria separando hospitales.
\end{column}

\begin{column}{0.48\textwidth}
\textbf{Region de interes}
\scriptsize

El campo irradiado se delimita como region, no como rectangulo: recortar deja
entre 20\% y 32\% de fondo porque el area irradiada no es rectangular.

No se rellenan huecos: las axilas se registran al nivel del fondo y
rellenarlas haria que la fraccion de area midiera el colimador.
\end{column}
\end{columns}
\end{frame}

% ---------------------------------------------------------------- 5
\begin{frame}{Resultado, tabla y reproduccion}
\footnotesize
\begin{columns}[T]
\begin{column}{0.5\textwidth}
\textbf{Ninguna estrategia de realce cumple el criterio}
\scriptsize

Criterio fijado antes de procesar: mejorar la linea base en 3 de 5
descriptores, en \textbf{ambos} origenes.

\vspace{0.15cm}
\begin{tabular}{@{}lcc@{}}
\toprule
Estrategia & MC & SH \\
\midrule
clahe\_fijo & 2/5 & 4/5 \\
eq\_adap & 3/5 & 1/5 \\
\textit{las otras cuatro} & \multicolumn{2}{c}{\textit{peor}} \\
\bottomrule
\end{tabular}

\vspace{0.2cm}
La mejor es \textbf{\MejorEstrategia{}}, sin realce: Dice \MejorDice{}
(\MejorDiceMC{} / \MejorDiceSH{}). Un promedio agregado habria hecho pasar a
clahe\_fijo.
\end{column}

\begin{column}{0.48\textwidth}
\textbf{Tabla y prevencion de fuga}
\scriptsize

\NFilasTabla{} filas $\times$ \NColumnasTabla{} columnas. Predictores: los
cinco descriptores sobre la mascara \emph{obtenida}.

Excluidos de \texttt{model\_input.json}: todo \texttt{ref\_}, \texttt{err\_},
Dice y Jaccard. Derivan de la referencia, que no existe en un caso nuevo.

\vspace{0.15cm}
11 comprobaciones automaticas sobre la tabla, todas superadas.

\vspace{0.25cm}
\textbf{Siguiente decision}
\scriptsize

Verificar si base\_adap mantiene su ventaja en validacion, y si la simetria
izq-der sigue siendo el descriptor mas inestable una vez separado el fallo de
segmentacion de la variacion anatomica.
\end{column}
\end{columns}
\end{frame}

\end{document}